\documentclass{resume} % Use the custom resume.cls style
\usepackage{fontawesome}
\usepackage{hyperref}
\usepackage{lipsum}
\usepackage{graphicx}
\usepackage{adjustbox}
\usepackage{subcaption}
\usepackage{setspace}
\usepackage[inline]{enumitem}
\usepackage{xcolor}
\usepackage[backend=biber,style=ieee,style=authoryear,maxbibnames=99,defernumbers=true]{biblatex}
\nocite{*}
\date{\today}
\addbibresource{pub.bib}
\makeatletter
\usepackage{hanging}
% This command ignores the optional argument for itemize and enumerate lists
\newcommand{\inlineitem}[1][]{%
\ifnum\enit@type=\tw@
    {\descriptionlabel{#1}}
  \hspace{\labelsep}%
\else
  \ifnum\enit@type=\z@
       \refstepcounter{\@listctr}\fi
    \quad\@itemlabel\hspace{\labelsep}%
\fi}
\makeatother
\parindent=0pt
\def\sectionskip{\smallskip} % Tighten the gap after each section heading (overrides resume.cls \medskip)
\setlength{\parskip}{0.45em} % Slightly tighten inter-paragraph/subsection spacing (parfill default is larger)
\usepackage[left=0.3in,top=0.3in,right=0.3in,bottom=0.3in]{geometry} % Document margins
\hyphenpenalty=10000
\DeclareNameAlias{sortname}{given-family}
\DeclareNameAlias{default}{given-family}
\DeclareNameFormat{author}{%
  \nameparts{#1}% split into given/family/etc.
  \ifboolexpr{
    test {\ifdefstring{\namepartfamily}{Datta}}
    and
    test {\ifdefstring{\namepartgiven}{Sohom}}
  }
    {\textbf{\namepartgiven}\space\textbf{\namepartfamily}}
    {\namepartgiven\space\namepartfamily}%
  \ifthenelse{\value{listcount}<\value{liststop}}
    {\addcomma\space}
    {}%
}
\AtEveryBibitem{%
  % Clear the default formatting
  \clearfield{url}
  \clearlist{language}
}

% Redefine inproceedings entries (conference papers)
\DeclareBibliographyDriver{inproceedings}{%
  % --- Title first ---
  \textbf{\printfield[titlecase]{title}}\\
  % --- Authors ---
  \printnames{author}\\
  % --- Venue ---
  \printfield{booktitle}%
  \iffieldundef{year}{}{\space(\printfield{year})}%
  \finentry
}

% Redefine article entries (journals)
\DeclareBibliographyDriver{article}{%
  \printfield[titlecase]{title}%
  \newunit\newblock
  \printnames{author}%
  \newunit\newblock
  \printfield{journaltitle}%
  \iffieldundef{year}{}{\space(\printfield{year})}%
  \finentry
}
\setlength\bibhang{0pt}
\defbibenvironment{bibliography}
  {\begin{list}{$\bullet$}{\leftmargin=1em \itemindent=0em \itemsep -0.6em}}
  {\end{list}}
  {\item}
\newcommand{\tab}[1]{\hspace{.\textwidth}\rlap{#1}}
\newcommand{\itab}[1]{\hspace{0em}\rlap{#1}}
\name{Sohom Datta} % Your name
\address{sohomdatta1+web(at)gmail.com}% Your phone number and email
\address {\faWikipediaW \href{https://gerrit.wikimedia.org/r/q/owner:sohomdatta1@gmail.com}{ MediaWiki Gerrit} \\
{\faGithub \href{https://github.com/sohomdatta1}{ GitHub}} \\
{\faLinkedin \href{https://www.linkedin.com/in/sohom-datta-c29ob20k/}{ Linkedin}}}%
\begin{document}
%\vspace{-1em}
\begin{rSection}{Education}
\begin{list}{$\bullet$}{\leftmargin=1em \itemindent=0em}
\itemsep -0.6em
\item {\bf North Carolina State University, Raleigh, USA} \hfill {\em August 2024 -- Present}  \\ Ph.D. in Computer Science. Advised by Dr. Alexandros Kapravelos (Focus: Web security and privacy)\hfill
\item {\bf Manipal Institute of Technology, Manipal, India} \hfill {\em July 2019 -- August 2023}  \\ B.Tech. in Computer Science and Engineering. (Coursework: Compilers, Software Testing, Computer Security)\hfill
\end{list}
\end{rSection}
\begin{rSection}{Research Interests}
Web security and privacy; browser internals and instrumentation; dynamic program analysis; large-scale measurement of online tracking; fuzzing and vulnerability discovery.
\end{rSection}
\rHeading{Publications}
\printbibliography[heading=none]
\begin{rSection}{Research Experience}
\begin{rSubsection}{\bf Graduate Research Assistant, Wolfpack Security and Privacy Research Lab}{\em January 2024 -- Present}{North Carolina State University}{}
\item Conduct research into privacy-violating JavaScript across heterogeneous execution environments, including mobile browsers and Android WebViews.
\item Maintain and extend VisibleV8, an open-source browser instrumentation framework, across 25+ major Chrome releases.
%\item Built the VisibleV8 Crawler, an open-source tool that enabled internet-scale (around 1M websites) web crawling with VisibleV8.
\item Ported VisibleV8 to Android WebViews, enabling the first dynamic analyses of privacy violations in the Android advertising ecosystem across 1K apps.
\end{rSubsection}
\begin{rSubsection}{\bf Visiting Researcher, Software Security Research Group}{\em September 2023 -- December 2023}{Max Planck Institute for Security and Privacy}{}
\item Studied race-condition bugs in application-layer logic across widely used web servers such as Express and Next.js.
\item Implemented fuzzers to automatically detect and flag security-critical race-condition issues in web server application logic.
\end{rSubsection}
\begin{rSubsection}{\bf Undergraduate Research Intern, Wolfpack Security and Privacy Research Lab}{\em February 2023 -- July 2023}{North Carolina State University}{}
%\item Worked on improving VisibleV8, a research tool aimed at making it easier to perform large-scale measurements of web security and abuse patterns on the internet.
\item Re-architected the VisibleV8 crawling pipeline, achieving a 20x improvement in crawling and post-processing of logs generated while visiting websites.
\item Contributed patches addressing deficiencies in VisibleV8's logging and tracing capabilities, including its ability to trace \texttt{eval(...)} and other code-execution pathways.
%\item Conducted research into measuring privacy data leakages across cross-party contexts using VisibleV8.
\end{rSubsection}
% \begin{rSubsection}{\bf Research Assistant, Information and Communication Dept.}{\em December 2021 -- June 2022}{Manipal Institute of Technology}{}
% \item Conducted research into inter-process communication mechanisms and their application as a data-exfiltration vector in cross-process Spectre attacks under Professor Nisha P Shetty.
% \item Worked on building a novel obfuscation technique aimed at reducing the efficacy of linkage attacks against ego-centric networks.
% \end{rSubsection}
\end{rSection}
\begin{rSection}{Honors and Awards}
\begin{list}{$\bullet$}{\leftmargin=1em \itemindent=0em}
\itemsep -0.5em
\item {\bf NC State University Graduate Fellowship } \hfill {\em 2024 -- 2025}
\item {\bf Google ESCAL8 2024 invitee } \hfill {\em 2024}
\item {\bf Google Vulnerability Rewards Program } -- Multiple awards for finding vulnerabilities across Go, Dart and Google Chrome \hfill {\em 2021 -- 2023}
\item {\bf Mozilla Hall of Fame } -- Multiple awards for finding vulnerabilities in Firefox \hfill {\em 2022 -- 2023}
\item {\bf CSAW CTF '21 (hosted by NYU) } -- placed among the top 15 teams in India; ASIS CTF Finals 2021 -- 2nd in India; ranked among the top 12 teams in India on CTFtime. \hfill {\em 2021--2022}
\end{list}
\end{rSection} 
% \item Google Vulnerability Rewards Program -- 7,500 USD for an XSS sanitization deficiency in the Go \texttt{html/template} library. \hfill {\em 2023}
% \item Google Chrome Vulnerability Rewards Program -- 5,000 USD for a mechanism to reliably leak a user's browsing history via an experimental ``origin-trial'' web feature in Chrome 116. \hfill {\em 2023}
% \item Google Vulnerability Rewards Program -- 3,133.70 USD for authentication bypasses in the \texttt{dart:core} URI parsing module in Dart. \hfill {\em 2022}
% \item Google Vulnerability Rewards Program -- 3,133.70 USD for a URL validation bypass in Google's Clojure library. \hfill {\em 2021}
% \item CSAW CTF '21 (hosted by NYU) -- placed among the top 15 teams in India; ASIS CTF Finals 2021 -- 2nd in India; ranked among the top 12 teams in India on CTFtime. \hfill {\em 2021--2022}
% \end{list}
% \end{rSection}
\begin{rSection}{Selected Vulnerability Disclosures}
\begin{list}{$\bullet$}{\leftmargin=1em \itemindent=0em}
\itemsep -0.5em
\item Authentication bypass, cross-site scripting, and information-disclosure vulnerabilities in English Wikipedia. (\texttt{CVE-2024-47848}, \texttt{CVE-2024-23174}, \texttt{CVE-2023-45369})
\item XSS sanitization deficiency in the Go \texttt{html/template} library. (\texttt{CVE-2023-24538})
\item Cross-site leakage (XS-Leak) issues in Google Chrome and Firefox's implementation of the ResourceTiming API. (\texttt{CVE-2022-1146}, \texttt{CVE-2022-29915})
\item Authentication bypasses in the \texttt{dart:core} URI parsing module in Dart. (\texttt{CVE-2022-3095}, \texttt{GHSA-m9pm-2598-57rj})
\item High-severity denial-of-service vulnerability in the jpeg-js JavaScript library. (\texttt{CVE-2022-25851}, \texttt{SNYK-JS-JPEGJS-2859218})
\end{list}
\end{rSection}
\begin{rSection}{Teaching, Mentoring, and Outreach}
\begin{rSubsection}{\bf Shadow Teaching Assistant}{\em 2024 -- 2026}{North Carolina State University}{}
\item Held office hours for CSC 405, CSC 537, and CSC 591, under Dr. Alexandros Kapravelos.
\end{rSubsection}
\begin{rSubsection}{\bf Mentor, HackPack}{\em 2024 -- 2025}{North Carolina State University}{}
\item Designed the infrastructure for and co-organized HackPackCTF, a capture-the-flag competition with 1000+ participants worldwide.
\item Mentored 7+ students in competing in CTFs, finding real-world security issues and working on academic research projects in web security and privacy.
\end{rSubsection}
\begin{rSubsection}{\bf Mentor, Google Summer of Code}{\em 2021 -- 2024}{Wikimedia Foundation}{}
\item Mentored 4 contributors through the Google Summer of Code program -- one as co-mentor (2021) and three as primary mentor (2023--2024).
\item Served as organization administrator for Wikimedia (2023) and attended the Google Mentor Summit (2024).
\end{rSubsection}
\begin{rSubsection}{\bf Member, Cryptonite Manipal (cybersecurity student group)}{\em April 2021 -- August 2022}{Manipal Institute of Technology}{}
\item Led a cybersecurity team of 22 students to top rankings among Indian teams in international CTFs.
\item Created challenges and managed the security infrastructure for niteCTF, an international CTF event that attracted 1200+ participants from 43 countries.
\item Conducted workshops on control-flow integrity, binary exploitation, and format-string exploits, covering state-of-the-art mitigations such as ASLR, stack canaries, and fuzz testing.
\end{rSubsection}
\end{rSection}
\begin{rSection}{Software and Open-Source Contributions}
\begin{rSubsection}{\bf Volunteer Maintainer, MediaWiki}{\em October 2021 -- Present}{Wikimedia Foundation}{}
\item Maintain the infrastructure for the ProofreadPage MediaWiki extension, which adds proofreading capabilities and is deployed across 70+ production wikis managed by Wikimedia. (80+ major patches)
\item Overhauled the preload and caching mechanism to reduce load times for editor-facing components.
\item Refactored and rewrote PageTriage, a tool used by volunteers to triage new pages created on Wikimedia wikis, to make it more modular and easier to maintain.
\item Was granted commit access across all projects under the MediaWiki umbrella in 2025.
% \item Built EditInSequence, a community-requested feature that allows users to edit multiple pages via a fast and easy to use interface.
\end{rSubsection}
\begin{rSubsection}{\bf Student Developer, Chromium}{\em June 2022 -- November 2022}{Google Summer of Code}{}
\item Aligned Chromium's implementation of the Performance API with the relevant W3C specifications.
\item Discovered and fixed 3 bugs related to incorrect First Contentful Paint and Largest Contentful Paint entries while refactoring how the PaintTiming API marked contentful images in Chromium.
\item Collaborated with the W3C Web Performance Working Group to rectify how LCP timing entries were reported at the painting layer.
\item Discovered and patched systemic cross-site leakage issues in the ResourceTiming API across the major browser engines WebKit, Blink, and Gecko (Firefox). (\texttt{CVE-2023-1232})
\end{rSubsection}
\begin{rSubsection}{\bf Student Developer, Wikimedia}{\em March 2020 -- September 2020}{Google Summer of Code}{}
\item Developed a feature that made it easier for volunteers to work with ``pagelists,'' a custom syntax for storing page-number information for multi-page files.
\item Developed heuristics that let users create and edit ``pagelists'' without interacting with the custom XML tag-based syntax.
\item Assisted in developing automated tests to detect accidental or malicious changes in minified JavaScript blobs imported as dependencies.
\end{rSubsection}
\vspace{-0.3em}
{\bf Selected research artifacts and projects}
\begin{list}{$\bullet$}{\leftmargin=1em \itemindent=0em}
\itemsep -0.5em
\item {\bf VisibleV8 crawler} -- an open-source tool enabling internet-scale ($\sim$1M websites) web crawling with VisibleV8, used extensively within the Wolfpack Security and Privacy Research Lab.
\item {\bf WebViewTracer} -- one of the first browser instrumentation frameworks for Android WebViews
\item {\bf Fuzzing sudo} -- ran a fuzzing campaign on sudo and found latent use-after-frees, out-of-bounds reads, and integer overflows. (sudo-project/sudo \texttt{Issue\#198}, \texttt{PR\#196}, \texttt{PR\#218}, \texttt{PR\#227})
% \item {\bf seccomp-bpf software sandbox} -- a process sandbox built on the kernel seccomp API that selectively allows or denies syscall sequences deemed malicious by a set of heuristic rules.
%\item {\bf Dis, a stack based language } - Developed a compiler in Typescript for a toy stack-based language that used polish notation to perform arithmetic operations and compiled into x86 assembly.
\end{list}
\end{rSection}
\begin{rSection}{Professional Service}
\begin{list}{}{\leftmargin=1em \itemindent=0em}
\itemsep -0.5em 
\item {\bf Artifact Reviewer}, IEEE Symposium on Security and Privacy (S\&P) 2026 Artifact Evaluation Committee \hfill {\em 2025 -- 2026}
\item {\bf Member}, Wikimedia Product and Technology Advisory Council \hfill {\em 2024 -- Present}
\item {\bf Member}, Wikimedia Unsupported Tools Working Group \hfill {\em 2025 -- Present}
%\item {\bf Shadow TA} for CSC 405, Computer Security North Carolina State University \hfill {\em Spring 2026}
\end{list}
\end{rSection}
\begin{rSection}{Technical Skills}
\begin{tabular}{ @{} >{\bfseries}l @{\hspace{6ex}} l }
Programming Languages \ & JavaScript, C, C++, Python, Go \\
Research Areas \ & Browser Internals, Web Security, Dynamic Analysis, Fuzzing \\
Tools \& Frameworks \ & Chromium, Puppeteer, Ghidra, IDA, AFL++, TensorFlow, PostgreSQL, MongoDB \\
\end{tabular}
\end{rSection}
\hrule
\small \small \begin{center}This resume was last updated on \DTMnow. The latest version of this resume is available \href{https://sohomdatta1.github.io/cv/artifacts/cv.pdf}{on Github}.\\
\color{white} The specific commit from which this resume was generated is .\end{center}
\end{document}